Une spire carrée de côté $a = \qty{\v{a}}{cm}$ et de résistance \qty{\v{R}}{\ohm} se déplace à vitesse constante $v = \qty{\v{v}}{m/s}$, dans son propre plan. Elle sort d'une région où règne un champ magnétique uniforme de \qty{\v{B}}{T}, perpendiculaire au plan de la spire. Le bord de la région est rectiligne et parallèle à un côté de la spire.

\begin{enumerate}[label=\alph*)]
    \item Pendant que la spire sort du champ, déterminez la force électromotrice induite.
    \item Déterminez le courant induit.
    \item Déterminez la force qu'il faut exercer sur la spire pour garder sa vitesse constante.
    \item Déterminez l'énergie dissipée dans la spire pendant toute sa sortie du champ.
\end{enumerate}
